\documentclass[12pt]{article}

% --- Paquets standards ---
\usepackage[french]{babel}
\usepackage[utf8]{inputenc}
\usepackage[T1]{fontenc}
\usepackage{geometry}
\geometry{a4paper, margin=2.5cm}

% --- Mise en page ---
\usepackage{setspace}
\onehalfspacing % Espacement pour faciliter la lecture orale

\title{Oral de Culture Générale : Consommation et Identité \\ \large Focus : Baudrillard, Bourdieu et le Jeu Vidéo}
\author{Maël - BTS SIO}
\date{Février 2026}

\begin{document}

\maketitle

\section*{Axe : Baudrillard et la Valeur-Signe dans le Virtuel}

Pour aborder la question de la consommation d'identités, il faut d'abord se pencher sur la pensée de \textbf{Jean Baudrillard}. Dans son analyse, il postule que nous avons dépassé le stade de l'utilité. 

\[ \text{Valeur-Signe} > \text{Valeur d'Usage} \]

Dans le cadre des jeux vidéo, cette théorie devient d'une clarté absolue. Prenons l'exemple d'un \textit{skin}. Techniquement, cet objet numérique n'a aucune valeur d'usage : il ne modifie pas les statistiques, il n'aide pas à gagner. C'est du pur code esthétique.

Ce que le joueur consomme, c'est donc un \textbf{signe}. On n'achète pas un objet, on achète ce qu'il communique aux autres. Baudrillard parle de \textbf{simulacre} : l'image (l'avatar) finit par devenir plus réelle socialement que l'individu réel. En achetant une apparence, le joueur n'équipe pas un personnage, il s'approprie une identité visuelle immédiatement reconnaissable dans un système de signes donné.

\clearpage

\section*{Axe : Bourdieu et la Distinction Sociale}

Si Baudrillard s'intéresse au signe, \textbf{Pierre Bourdieu} nous permet de comprendre comment ces signes servent à hiérarchiser les individus. Pour Bourdieu, la consommation est une arme de \textbf{distinction}.

L'acte d'achat répond à une logique de positionnement social :
\[ \text{Habitus} + \text{Capital (Économique/Culturel)} = \text{Pratiques de Distinction} \]

Dans l'univers vidéoludique, cette distinction se traduit par le prestige lié à la rareté. Posséder un skin exclusif ou ancien n'est pas qu'une question d'esthétique, c'est l'affichage d'un \textbf{capital symbolique}. Cela prouve que le joueur possède les moyens financiers ou l'ancienneté (le code) pour appartenir à l'élite du jeu.

À l'inverse, Bourdieu évoque la \textbf{violence symbolique}, que l'on retrouve dans le mépris envers les « no-skins » ou « bambis ». Celui qui ne consomme pas est exclu ou stigmatisé. Ici, la consommation d'identité est une stratégie de survie sociale : on consomme pour ne pas être classé parmi les « inférieurs ». L'identité devient ainsi un marqueur de classe au sein d'une société virtuelle structurée.

\end{document}